% GENERATED FILE. Edit canonical lessons, then run npm run generate:books.
% canonical-source-hash: fnv1a64:f7fca58a71d9e4ff
% canonical-lessons: LA-C250-vitis, LA-C250-quercus, LA-C250-pinus, LA-C250-laurus, LA-C250-viola

\chapter{Vine, Oak, Pine, Laurel, Violet}
\label{ch:la-vine-oak-pine-laurel-violet}
\hlchaptermodality{250}

\begin{chapteropening}
\textbf{By the end of this chapter:} \emph{I can say a vine, an oak, a pine, a laurel and a violet.}

Complete the last lesson of chapter 250: I can say a vine, an oak, a pine, a laurel and a violet.
\end{chapteropening}

\section[vītis, vītis]{vītis, vītis — a vine}

\label{lesson:LA-C250-vitis}



\begin{quote}
Before the new one: say the Latin for hide, then the Latin for linen.
\end{quote}

\subsection*{You'll want to know: vītis, vītis}
\textbf{vītis, vītis} — \textquotedblleft{}a vine\textquotedblright{}.

\begin{cousinweb}
\textbf{vītis} is inherited, usually traced to Proto-Indo-European \emph{*weh1itis}, \textquotedblleft{}a twining plant\textquotedblright{}, from \emph{*weh1y-}, \textquotedblleft{}to twist, to weave\textquotedblright{} (EtymDB 2.1); \emph{vieō}, \textquotedblleft{}to plait\textquotedblright{}, is from the same root (Lewis \& Short).
\end{cousinweb}

\begin{grammarlens}[title={What you've built}]
One more word for this chapter.
\end{grammarlens}

\subsection*{Your turn}
\begin{itemize}
\raggedright
  \item \emph{Say it:} \emph{vītis, vītis}
  \item \emph{Say it:} \emph{vītis, vītis}, once more
  \item \emph{Say it:} say \emph{vītis, vītis}
  \item \emph{Recall:} say the Latin for hide, then the Latin for linen, then say \emph{vītis, vītis} again
\end{itemize}

\begin{tcolorbox}[breakable,colback=teal!4,colframe=teal!35!black,title={Before you move on}]
What does vītis, vītis mean? (\textquotedblleft{}A vine\textquotedblright{}.) Say it once more.
\end{tcolorbox}
\section[quercus, quercūs]{quercus, quercūs — an oak}

\label{lesson:LA-C250-quercus}



\begin{quote}
Before the new one: say the Latin for linen, then the Latin for a vine.
\end{quote}

\subsection*{You'll want to know: quercus, quercūs}
\textbf{quercus, quercūs} — \textquotedblleft{}an oak\textquotedblright{}.

\begin{cousinweb}
\textbf{quercus} is inherited, probably from Proto-Italic \emph{*kwerkus}, \textquotedblleft{}oak\textquotedblright{} (EtymDB 2.1). It is usually traced to Proto-Indo-European \emph{*perkwus}, \textquotedblleft{}oak\textquotedblright{}, with the first sound changed to match the second, as in \emph{coquō}, \textquotedblleft{}to cook\textquotedblright{}. English \emph{fir} is a distant cousin.
\end{cousinweb}

\begin{grammarlens}[title={What you've built}]
One more word for this chapter.
\end{grammarlens}

\subsection*{Your turn}
\begin{itemize}
\raggedright
  \item \emph{Say it:} \emph{quercus, quercūs}
  \item \emph{Say it:} \emph{quercus, quercūs}, once more
  \item \emph{Say it:} say \emph{quercus, quercūs}
  \item \emph{Recall:} say the Latin for linen, then the Latin for a vine, then say \emph{quercus, quercūs} again
\end{itemize}

\begin{tcolorbox}[breakable,colback=teal!4,colframe=teal!35!black,title={Before you move on}]
What does quercus, quercūs mean? (\textquotedblleft{}An oak\textquotedblright{}.) Say it once more.
\end{tcolorbox}
\section[pīnus, pīnūs]{pīnus, pīnūs — a pine}

\label{lesson:LA-C250-pinus}



\begin{quote}
Before the new one: say the Latin for a vine, then the Latin for an oak.
\end{quote}

\subsection*{You'll want to know: pīnus, pīnūs}
\textbf{pīnus, pīnūs} — \textquotedblleft{}a pine\textquotedblright{}.

\begin{cousinweb}
\textbf{pīnus} is inherited, usually traced to Proto-Indo-European \emph{*piHnu-}, \textquotedblleft{}pine\textquotedblright{}, from \emph{*peyH-}, \textquotedblleft{}sap, resin\textquotedblright{} (EtymDB 2.1); Greek \emph{pitys}, \textquotedblleft{}pine\textquotedblright{}, is a cousin (Lewis \& Short).
\end{cousinweb}

\begin{grammarlens}[title={What you've built}]
One more word for this chapter.
\end{grammarlens}

\subsection*{Your turn}
\begin{itemize}
\raggedright
  \item \emph{Say it:} \emph{pīnus, pīnūs}
  \item \emph{Say it:} \emph{pīnus, pīnūs}, once more
  \item \emph{Say it:} say \emph{pīnus, pīnūs}
  \item \emph{Recall:} say the Latin for a vine, then the Latin for an oak, then say \emph{pīnus, pīnūs} again
\end{itemize}

\begin{tcolorbox}[breakable,colback=teal!4,colframe=teal!35!black,title={Before you move on}]
What does pīnus, pīnūs mean? (\textquotedblleft{}A pine\textquotedblright{}.) Say it once more.
\end{tcolorbox}
\section[laurus, laurī]{laurus, laurī — a laurel}

\label{lesson:LA-C250-laurus}



\begin{quote}
Before the new one: say the Latin for an oak, then the Latin for a pine.
\end{quote}

\subsection*{You'll want to know: laurus, laurī}
\textbf{laurus, laurī} — \textquotedblleft{}a laurel\textquotedblright{}.

\begin{cousinweb}
\textbf{laurus} has no sure origin; it is usually thought to be a loan from a lost Mediterranean language. German \emph{Lorbeer}, \textquotedblleft{}laurel\textquotedblright{}, was borrowed from it.
\end{cousinweb}

\begin{grammarlens}[title={What you've built}]
One more word for this chapter.
\end{grammarlens}

\subsection*{Your turn}
\begin{itemize}
\raggedright
  \item \emph{Say it:} \emph{laurus, laurī}
  \item \emph{Say it:} \emph{laurus, laurī}, once more
  \item \emph{Say it:} say \emph{laurus, laurī}
  \item \emph{Recall:} say the Latin for an oak, then the Latin for a pine, then say \emph{laurus, laurī} again
\end{itemize}

\begin{tcolorbox}[breakable,colback=teal!4,colframe=teal!35!black,title={Before you move on}]
What does laurus, laurī mean? (\textquotedblleft{}A laurel\textquotedblright{}.) Say it once more.
\end{tcolorbox}
\section[viola, violae]{viola, violae — a violet}

\label{lesson:LA-C250-viola}



\begin{quote}
Before the new one: say the Latin for a pine, then the Latin for a laurel.
\end{quote}

\subsection*{You'll want to know: viola, violae}
\textbf{viola, violae} — \textquotedblleft{}a violet\textquotedblright{}.

\begin{cousinweb}
\textbf{viola} is related to Greek \emph{ion}, \textquotedblleft{}violet\textquotedblright{} (Lewis \& Short), which once began with a w sound; both are probably loans from a lost Mediterranean language.
\end{cousinweb}

\begin{grammarlens}[title={What you've built}]
One more word for this chapter.
\end{grammarlens}

\subsection*{Your turn}
\begin{itemize}
\raggedright
  \item \emph{Say it:} \emph{viola, violae}
  \item \emph{Say it:} \emph{viola, violae}, once more
  \item \emph{Say it:} say \emph{viola, violae}
  \item \emph{Recall:} say the Latin for a pine, then the Latin for a laurel, then say \emph{viola, violae} again
\end{itemize}

\begin{tcolorbox}[breakable,colback=teal!4,colframe=teal!35!black,title={Before you move on}]
What does viola, violae mean? (\textquotedblleft{}A violet\textquotedblright{}.) Say it once more.
\end{tcolorbox}
